\documentclass[../../main-detail.tex]{subfiles}
\begin{document}
\chapter{翻訳集}
\section{りんご文}
\subsection{ざすろんリスト}
\todo{5-2からの移植．あとで現文法に合わせて書き換える．}
\begin{exe}
    \ex \begin{xlist}
        \ex kaja hamin ke kalam
        \ex kaja cu hamin ke kalam kon mefante
        \ex (傾向構文)
        \todo{傾向構文の考察がなされたら作文する}
    \end{xlist}
    \glt 私はりんごを食べる．
    \glt \emph{解説}：a.は時点を問わず(未来もあり得る) 「私がりんごを食べる」行為を宣言している．b.は私が少なくない数の「りんごを食べる」行為をすることを表している．c.は「私はりんごを食べる人だ」ということを表している．
\end{exe}
\begin{exe}
    \ex kaja haminik ke kalam
    \glt 私はりんごを食べた．
    \ex toja haminen ke kalam
    \glt 彼はりんごを食べている．
    \ex toja meihaen hamin ke kalam
    \glt 彼女はりんごを食べ終わっている．
    \ex toja meihaik hamin ke kalam
    \glt 彼女はりんごを食べ終わっていた．
    \ex hotaltoalka ja hamin kon ke kalam kon
    \glt 私の妻はりんごを食べたことがある．
    \ex kostafe hotaltoalka ja hamin ke kalam aki nos san
    \glt 私の妻はりんごを毎日食べる．
    \glt konstafe $\varphi$ (tos)によってそれ以降の$^*$を全て$\exists$で量化することができる．
    \ex \begin{xlist}
        \ex mika hotaltoalka san ja hamin ke kalam kon aki senos
        \ex kostafe mika hotaltoalka san ja hamin ke kalam aki senos
    \end{xlist}
    \glt 私と私の妻は昨日りんごを食べた．
    \glt a.は一緒に食べた場合(食べたりんごは別)，b.は別々に食べた場合である．
    \ex mika hotaltoalka san ja hamin ke kalam kon aki noin senos
    \glt 私と私の妻は6日前にりんごを食べた．
    \ex cotole san ja hamin ke kalam kon aki lonos
    \glt 彼らは明日りんごを食べる．
    \ex cotole san ja hamin ke kalam kon aki noin lonos
    \glt 彼らは6日後にりんごを食べる．
    \ex cu kostafe cotole ja hamin ke leta aki nosle lulas nais sas kalam
    \glt 彼女らは3日間りんごを食べている．（食事としてりんごだけを食べている）
    \glt 直訳すると「その3日間で彼らが食べたものは全てりんごである」となる．
    \ex kostafe cotole san hamin ke kalam aki selovfon hanijo san cistoa la taksal
    \glt 彼女らはりんごを5分間食べ続けている．
    \ex toja haminlefla ke kalam kon aki selov san
    \glt 彼は常にりんごを食べている．（四六時中ずっとりんごを食べている）
    \ex lulas kalamle
    \glt りんごが3つある．
    \ex kalamkla tov sastov stanokal
    \glt りんご達がテーブルの上にある．
    \ex coto kon ja haminik ke kalam
    \glt 誰かがあのりんごを食べてしまった．（その結果，あのりんごは今はない）
    \ex nockom to min coto kon ja haminas ke kalam
    \glt 誰かがこのりんごを食べそうだ．（推量）
    \ex nockom to min coto kon ja haminas ke kalam mi kalam mi kalam san
    \glt 誰かがこのりんご，そのりんご，あのりんごを食べたそうだ．（伝聞）
    \ex linlas fe leta sas kalam
    \glt りんごがひとつもない．
    \glt leta sas kalamはkalamの冪集合であり，feでその中から1つ選んでいることを明示している．
    \ex \begin{xlist}
        \ex meide kaja hamin ke kalam
        \ex meide kostafe kaja hamin ke kalam
    \end{xlist}
    \glt 私はりんごを食べない．
    \glt a.は特定のりんごを，b.はいかなるりんごをも食べないという意味になる．
    \ex meide kaeli to min kostafe kaja hamin ke kalam
    \glt 私はりんごを食べられない．（りんごがないので状況的に食べられない）
    \ex kaja khisuik e kalam na noolja lja saida
    \glt 私はりんごを床に落とした．
    \glt 直訳すると「私はりんごを落として(それは)床に行った」となる．
    \ex khisuik e kalam na noolja lja saida
    \glt りんごが床に落ちた．
    \glt 直訳すると「りんごを落として(それは)床に行った」となる．能動態/受動態の区別はない(意図的にそう設計している)．
    \ex kaeli beto min kalam kon ke hamin kon
    \glt りんごは食べられない．（りんごの性質として，可食でなくなった）
    \ex \begin{xlist}
        \ex tomja to cu kaja le hamin ke kalam nais anic
        \ex tomja to min kaja hamin ke kalam
    \end{xlist}
    \glt 私はりんごを食べたい．
    \glt a.はastavを使った書き方，b.はmintosを使った書き方である．
    \ex tomja to cu kaja le cutov e kalam nais anic
    \glt 私はりんごを買いたい．
    \ex zoniloc noccea kalam
    \glt このりんごは汚い．
    \ex meide kostafe lawnloc noccea kalam
    \glt このりんごは綺麗ではない．
    \ex nackom to min haminas ke kalam
    \glt あのりんごは食べられそうだ．
    \glt 直訳すると「あのりんごは食べられるかもしれない」となる
    \ex meide nato litoi kalam
    \glt あれはりんごではない．
    \ex nato belitoi kalam litoi mikam
    \glt これはりんごではなくみかんだ．
    \ex hatast noja cu le hamin ke kalam kon nais mefante
    \glt あなたはみかんを食べるか？
    \ex ai kaja cu le hamin ke kalam kon nais mefante
    \glt はい，私はみかんを食べます．
    \ex hatast meide noja cu le hamin ke kalam kon nais mefante
    \glt あなたはみかんを食べないのですか？
    \ex ai meide kaja cu le hamin ke kalam kon nais mefante
    \glt ええ，私はみかんを食べません．
    \ex hatast walno tomav to cu noja le hamin ke mikam kon nais anic no
    \glt あなたはみかんを食べたかった．ですよね？
    \glt 付加疑問文は未実装
    \ex mei meide walka tomav to cu kaja le hamin ke mikam kon nais anic no
    \glt いいえ，私はみかんを食べたくはありません．
    \ex hatast meide noja hamin ke kalam mi mikam kon
    \glt あなたはりんごもみかんも食べないのですか？
    \ex mei kaja hamin ke kalam
    \glt いや，私はりんごを食べます．
    \ex meide kaja hamin ke mikam tos suinot meide hamin ke mikam
    \glt 私はみかんを食べない．だから，みかんは食べられない．（受け身の否定）
    \glt 他の人が食べる可能性があるため，推論としては正しくないことに注意．
    \ex meide kaeli to min hamin ke kalam tos meide kostafe lawnloc noccea kalam
    \glt りんごは食べることができない．なぜなら，りんごは綺麗ではないから．
    \glt 否定文にsuinotを付けられない不具合がある．今後修正する．
    \ex zoniloc coi nocceanov cal kalam kojannot nocceanov coi lawn cal mikam
    \glt りんごは汚い．一方で，みかんは綺麗だ．
    \ex noja haminas ke mikam no
    \glt みかんを食べましょうよ？（勧誘）
    \ex mei
    \glt いいえ，お断りします．
    \ex hatast noja haminas ke mikam meide haminas ke kalam jo
    \glt あそこでみかんを食べろ．りんごは食べるな．
    \ex hatast tactno hactoskonazas ljaka sowano no tactka hactoskonazas ljano sowaka
    \glt りんごを頂いても宜しいですか？みかんは差し上げます．
    \ex kaja saikik me minac hatast meide noja haminas ke kalam noja haminas ke mikam no meeljaja
    \glt 私は「りんごは食べるな，みかんを食べろ」と言った．
    \ex hatast walno tomav to cu kaja le hamin ke kalam kon tos anic no
    \glt あなたはりんごを食べろと言うのですか？
    \glt 直訳すると「あなたは私にりんごを食べて欲しいの？」となる．
    \ex noja saikik me min walno tomav to cu noja le hamin ke kalam nais anic takhtal
    \glt 貴様はりんごを食べたいと言った．
    \ex meide walka lekonoy to min kaja saikik to min tomja to min kaja hamin ke kalam
    \glt 私はりんごを食べたいと言った覚えがない．
    \ex walka nackom to min noja saikik to minac tomja to cu le hamin ke kalam nais anic meeljaja
    \glt あなたは確か「私はりんごを食べたい」と言ったはずだ．
    \ex kaeli to min kaja saikik to saki
    \glt 私はもしかしたらそう言ったかもしれない．
    \ex walka kastav to min noja saikik min tomav to min hamin ke kalam takhtal na lagoley ljaka
    \glt あなたは絶対に私に対して「りんごを食べたい」と言った．
    \glt 直訳すると「あなたは『りんごを食べたい』と言って私に伝えたと私は信じている」となる．「絶対に」みたいな表現はまだ実装されていない．
    \ex tactato tovlent luis hamin kaja ke kalam
    \glt 彼は私にりんごを食べさせる．
    \ex meide kaja hamin ke mikam
    \glt みかんは私に食べられていない．
    \ex tactato tovlent luis hamin kaja ke kalam
    \glt りんごは彼によって私に食べさせられた．
    \ex meide konta tov kant
    \glt ここには何もない．
\end{exe}
\end{document}
